\section{CoHA to $\mathcal{W}_{1+\infty}$ to chiral Yangian to gauge theory: a rigorous treatise through three examples}
\label{wn:sec:coha-to-W-infty-treatise}

\providecommand{\ClaimStatusTheorem}{\textsc{[theorem, cited]}}
\providecommand{\ClaimStatusConjectured}{\textsc{[conjectural]}}
\providecommand{\ClaimStatusDefinition}{\textsc{[definition]}}
\providecommand{\ClaimStatusDerivation}{\textsc{[first-principles derivable]}}
\providecommand{\ClaimStatusOpen}{\textsc{[open]}}
\providecommand{\ClaimStatusCorrected}{\textsc{[corrected]}}
\providecommand{\ClaimStatusRetracted}{\textsc{[retracted]}}
\providecommand{\CoHA}{\mathrm{CoHA}}
\providecommand{\Sh}{\mathrm{Sh}}
\providecommand{\Hilb}{\mathrm{Hilb}}
\providecommand{\Wilpha}{\mathcal{W}_{1+\infty}}
\providecommand{\ChirY}{\mathcal{Y}^{\mathrm{ch}}}
\providecommand{\hCS}{\mathrm{hCS}}
\providecommand{\Conf}{\mathrm{Conf}}
\providecommand{\gghone}{\widehat{\mathfrak{gl}}_1}
\providecommand{\Fact}{\mathrm{Fact}}

This treatise develops three worked examples of the CoHA $\to$ vertex algebra $\to$
gauge theory chain, at successively increasing difficulty:
\begin{enumerate}
\item $\mathbb{C}^3$ (local, toric, flat): the Jordan triple loop quiver, affine Yangian
of $\mathfrak{gl}_1$, $\Wilpha$, and 4D $\mathcal{N} = 2$ $U(1)$ with $\Omega$-background.
\item Resolved conifold (local, toric, non-flat): the conifold quiver, CoHA, chiral
Yangian from holomorphic Chern--Simons.
\item $K3 \times E$ (compact, non-toric): obstructions to CoHA construction, partial
localisation strategies, Davison BPS Lie algebra, and the unresolved passage to
a chiral BKM.
\end{enumerate}

Every claim is tagged with scope.

%--------------------------------------------------------------------------------
\subsection{Example 1: $\mathbb{C}^3$ via the Jordan triple loop quiver}
\label{wn:subsec:jordan-triple}
%--------------------------------------------------------------------------------

\subsubsection{The quiver and its Jacobi algebra}

\ClaimStatusDefinition\ \emph{Jordan triple loop quiver $Q$.}
One vertex $*$; three loops $X, Y, Z: * \to *$.
Path algebra $kQ = k\langle X, Y, Z\rangle$ (free associative algebra on three generators).

\ClaimStatusDefinition\ \emph{Potential $W \in kQ_{\mathrm{cyc}} / [kQ, kQ]$.}
$W := \mathrm{tr}(X[Y, Z]) = \mathrm{tr}(XYZ - XZY)$.

\ClaimStatusDerivation\ \emph{Jacobi algebra.}
$J(Q, W) := kQ / \langle \partial_a W : a \in \{X, Y, Z\}\rangle$.
Compute the cyclic derivatives:
\begin{align}
\partial_X W &= YZ - ZY = [Y, Z],\\
\partial_Y W &= ZX - XZ = [Z, X],\\
\partial_Z W &= XY - YX = [X, Y].
\end{align}
Hence $J(Q, W) = k\langle X, Y, Z\rangle/([X, Y], [Y, Z], [Z, X]) = k[X, Y, Z]$.

\emph{Conclusion.} $J(Q, W) \cong \mathcal{O}_{\mathbb{C}^3}$, the coordinate ring of
affine three-space. The Jordan triple loop quiver with potential $\mathrm{tr}(XYZ - XZY)$
is a presentation of $\mathbb{C}^3$ as a Calabi--Yau three-fold. (Kontsevich--Soibelman
2008 \texttt{arXiv:0811.2435} §8; Ginzburg 2006 for CY3 algebras.)

\subsubsection{Moduli of representations and their equivariant cohomology}

\ClaimStatusDefinition\ A representation of dimension $n$ is a triple
$(X, Y, Z) \in \mathrm{End}(\mathbb{C}^n)^3$ satisfying $[X, Y] = [Y, Z] = [Z, X] = 0$
(commuting triple of $n \times n$ matrices).

\ClaimStatusDefinition\ \emph{Moduli stack.}
$\mathcal{M}_n := [\{(X, Y, Z) \in \mathrm{End}(\mathbb{C}^n)^3 : [X,Y]=[Y,Z]=[Z,X]=0\}/GL_n(\mathbb{C})]$.

\ClaimStatusDefinition\ \emph{Torus action.}
$T = (\mathbb{C}^\times)^3$ acts on $(X, Y, Z)$ with weights $(t_1, t_2, t_3)$.
The CY$_3$ condition (coming from the holomorphic volume form $dX \wedge dY \wedge dZ$
being $T$-invariant) is $t_1 t_2 t_3 = 1$, i.e.
$\epsilon_1 + \epsilon_2 + \epsilon_3 = 0$ where $\epsilon_i = \log t_i$.
In the equivariant parameter ring, we work over
$\mathbb{F} := \mathbb{C}(\epsilon_1, \epsilon_2)$ with $\epsilon_3 = -\epsilon_1 - \epsilon_2$.

\ClaimStatusDerivation\ \emph{Torus-fixed loci.}
A $T$-fixed commuting triple is necessarily diagonal: $X, Y, Z$ diagonalise
simultaneously with common eigenbasis. Fixed points are parameterised by
multisets of $n$ points in the weight lattice $\mathbb{Z}^3$. (This is the
standard fact that commuting triples of matrices decompose into weight spaces.)

For our purposes the relevant subcollection is the \emph{Hilbert scheme
stratum}: configurations where the eigenvalues are distinct and $T$-fixed.
These correspond to 3d partitions / plane partitions of size $n$.

\subsubsection{The CoHA multiplication}

\ClaimStatusDefinition\ (Kontsevich--Soibelman 2008 §6.)
\[
\CoHA(\mathbb{C}^3) := \bigoplus_{n \geq 0} H^*_T(\mathcal{M}_n, \phi_{W})
\]
where $\phi_W$ is the vanishing cycle sheaf of the potential $W$.

For the Jordan triple loop quiver with potential as above, the vanishing cycle
localises onto the commuting locus and the equivariant cohomology reduces to
localisation on $T$-fixed configurations. Over $\mathbb{F}$:
\ClaimStatusTheorem\ (Schiffmann--Vasserot 2013 \texttt{arXiv:1202.2756} Thm 1.1)
\[
\CoHA(\mathbb{C}^3) \otimes_\mathbb{F} \mathbb{F}(\!(z)\!) \cong \Sh
\]
where $\Sh$ is the shuffle algebra defined in §\ref{wn:sec:yplus-to-chiral-yangian-treatise}.

\ClaimStatusDefinition\ \emph{Shuffle multiplication, explicit form.}
For $F \in \mathbb{F}[z_1, \ldots, z_m]^{S_m}$ and $G \in \mathbb{F}[w_1, \ldots, w_n]^{S_n}$:
\[
(F \star G)(x_1, \ldots, x_{m+n}) :=
\frac{1}{m! n!} \sum_{\sigma \in S_{m+n}} \sigma \cdot \!\Big[
F(x_1, \ldots, x_m)\, G(x_{m+1}, \ldots, x_{m+n})
\prod_{\substack{1 \leq i \leq m \\ m < j \leq m+n}} \omega(x_j, x_i) \Big]
\]
with
\[
\omega(z, w) := \frac{(z - w - \epsilon_1)(z - w - \epsilon_2)(z - w - \epsilon_3)}{(z - w)^3}.
\]

\ClaimStatusDerivation\ \emph{Unit and gradation.}
The unit is $1 \in \Lambda_0 = \mathbb{F}$. Grading: $\deg(F) = m$ for $F \in \Lambda_m$.

\subsubsection{The positive half $Y^+(\gghone)$}

\ClaimStatusDefinition\ $Y^+(\gghone) := \CoHA(\mathbb{C}^3)_+$ := the sub-shuffle algebra
on symmetric polynomials with no poles as any $z_i \to z_j$ ($i \neq j$). (The full
shuffle algebra $\Sh$ allows poles that are killed by the $(z - w)^3$ denominator.)

\emph{This is the object traditionally called the positive half of the affine Yangian
of $\mathfrak{gl}_1$.} Under the Schiffmann--Vasserot isomorphism, it is generated by
elementary power sums $p_k := \sum_i z_i^k$ for $k \geq 0$.

\ClaimStatusDerivation\ \emph{Low-order multiplication.}
Using the prefactor $\tfrac{1}{m! n!}$ of the displayed shuffle formula: at $m = n = 1$
this prefactor equals $1$, and the $S_2$ sum expands to both orderings of $(z_1, z_2)$.
For $p_1 \star p_1$ with $p_1(z) = z$:
$(p_1 \star p_1)(z_1, z_2) = z_1 \cdot z_2 \cdot \omega(z_2, z_1)
+ z_2 \cdot z_1 \cdot \omega(z_1, z_2)
= z_1 z_2 [\omega(z_2, z_1) + \omega(z_1, z_2)]$.
(The apparent $\tfrac{1}{2}$ appearing in some normalisations compensates a double-counting
convention on the symmetric group; here the $\tfrac{1}{m!n!}$ convention already accounts
for this, so no extra $\tfrac{1}{2}$ arises.)

Expanding $\omega(z, w) + \omega(w, z)$ in large-$|z - w|$:
this is a symmetric rational function of $z - w$ which, after expansion, gives a
polynomial in $\epsilon_1, \epsilon_2, \epsilon_3$ times symmetric functions of $z, w$.

The explicit generators-and-relations presentation (Tsymbaliuk 2017
\texttt{arXiv:1703.04551}) gives $p_k$ satisfying
\[
[p_k, p_\ell] = \begin{cases}
\text{explicit polynomial in lower } p_j \text{ with coefficients in } \mathbb{F}, & k, \ell \geq 0.
\end{cases}
\]

\subsubsection{Drinfeld double and identification with $\mathcal{W}_{1+\infty}$-affine Yangian}

\ClaimStatusTheorem\ (Tsymbaliuk 2017 \texttt{arXiv:1703.04551} Thm~1.1.)
The Drinfeld double
\[
Y(\gghone) := Y^+ \bowtie Y^0 \bowtie Y^-
\]
is isomorphic, as a topological Hopf algebra, to the \emph{affine Yangian of
$\mathfrak{gl}_1$} in the Drinfeld-currents presentation with structure function
$\varphi(u) = (u + \epsilon_1)(u + \epsilon_2)(u + \epsilon_3)/u^3$.

\subsubsection{$\Wilpha$: two free-field definitions}

There are two routes to $\Wilpha$; their equivalence is non-trivial.

\ClaimStatusDefinition\ \emph{Kac--Radul 1993 definition.}
$\Wilpha[\lambda]$ is the unique one-parameter family of vertex operator algebras with
central charge $c = \lambda$, generated by a single field at each conformal dimension
$h = 1, 2, 3, \ldots$, with OPE structure determined by a $W$-infinity parent symmetry.

\ClaimStatusDefinition\ \emph{Free-field construction.}
$\Wilpha[N]$ is realised as the algebra of $\widehat{U(1)}^N$-invariant differential polynomials
in $N$ free bosons, at specific level. In the large-$N$ limit with appropriate
renormalisation this extends to a one-parameter family $\Wilpha[\lambda]$.

\ClaimStatusTheorem\ (Linshaw 2011 \emph{Invariant theory and $\Wilpha$}, also the
book by Linshaw \emph{Universal W-algebras}.)
The two definitions coincide: $\Wilpha[\lambda]$ is uniquely characterised by its
central charge and generating-field count.

\ClaimStatusTheorem\ (Schiffmann--Vasserot 2013 + Gaiotto--Rapčák 2017
\texttt{arXiv:1703.00982}.) The affine Yangian $Y(\gghone)$ maps to
$\Wilpha[\lambda]$-endomorphisms via the state-field correspondence on the
vacuum module.

\emph{Convention clarification.} Two $\lambda$-parameters coexist in the literature and
must not be conflated:
\begin{itemize}
\item $\lambda_{\mathrm{Tr}}$, the \emph{truncation parameter} in the SV / Tsymbaliuk
triangular-presentation convention, built from the pair $(\epsilon_1, \epsilon_2)$
before the CY$_3$ constraint:
\[
\lambda_{\mathrm{Tr}} := \frac{\epsilon_1}{\epsilon_3} + \frac{\epsilon_2}{\epsilon_3}
= \frac{\epsilon_1 + \epsilon_2}{\epsilon_3} \quad
\text{(with $\epsilon_3$ \emph{unconstrained})}.
\]
\item $\lambda_{\mathrm{GR}}$, the Gaiotto--Rap\v c\'ak parametrisation in which
the central charge of $\Wilpha[\lambda_{\mathrm{GR}}]$ is written as an affine-linear
functional of $(\epsilon_1, \epsilon_2, \epsilon_3)$ with the three $\epsilon_i$ on
symmetric footing.
\end{itemize}
The two are related by a linear-fractional transformation; neither is ``the'' $\lambda$.

\emph{CY$_3$: truncation collapse vs $\Wilpha$ family.} The constraint
$\epsilon_1+\epsilon_2+\epsilon_3=0$ forces
\[
\lambda_{\mathrm{Tr}}
=\frac{\epsilon_1+\epsilon_2}{\epsilon_3}
=-1.
\]
Thus the SV/Tsymbaliuk truncation parameter collapses on the CY$_3$ slice.
This does \emph{not} collapse the $\Wilpha$ family: the free parameter is
the Gaiotto--Rap\v c\'ak/$\Wilpha$ parameter
$\lambda_{\mathrm{GR}}$ (equivalently $\lambda_W$), which is related to
$\lambda_{\mathrm{Tr}}$ only after a convention-dependent
linear-fractional transformation and normalisation of the central charge.

\subsubsection{Gauge-theory origin: 4D $\mathcal{N} = 2$ $U(1)$ with $\Omega$-background}

\ClaimStatusDefinition\ \emph{Setup.} 4D $\mathcal{N} = 2$ pure $U(1)$ gauge theory on
$\mathbb{R}^4 = \mathbb{C}^2_{\epsilon_1, \epsilon_2}$ where the $\Omega$-background
rotates the two complex planes with weights $\epsilon_1, \epsilon_2$.

\ClaimStatusTheorem\ (Nekrasov 2003 \texttt{arXiv:hep-th/0206161}.)
The equivariant instanton partition function is
\[
Z^{\mathrm{inst}}_{\Omega}(\epsilon_1, \epsilon_2; q) = \sum_{n \geq 0} q^n \chi_T(\Hilb^n(\mathbb{C}^2))
= \prod_{n \geq 1} \frac{1}{1 - q^n}
\]
in the Abelian case, a free-boson-like generating function.

\ClaimStatusDefinition\ \emph{Holomorphic twist on spectral line.}
Add a further circle direction $\mathbb{C}_z$ (``spectral line'') so that the
total spacetime is $\mathbb{C}^2 \times \mathbb{C}_z$. The theory is topological in the
$\mathbb{C}^2$ directions (under $\Omega$-twist) and holomorphic in $\mathbb{C}_z$ (under
HT twist).

\ClaimStatusTheorem\ (Costello 2013 \texttt{arXiv:1303.2632} for the 4D
$\mathcal{N} = 1$ holomorphic-twist construction; upgraded to 4D $\mathcal{N} = 2$ with
$\Omega$-background in Costello--Gaiotto 2018 and Costello--Paquette 2020
\texttt{arXiv:2009.04834}; the general factorisation-algebra machinery is
Costello--Gwilliam 2017 \emph{Factorization Algebras in Quantum Field Theory} Vol II §5.)
The observable algebra of this twisted theory is a factorisation algebra
$\mathcal{F}_{4D \mathcal{N} = 2}^{U(1), \Omega}$ on $\mathbb{C}_z$ with holomorphic
locality (OPE singularities are poles in $z - w$).

\ClaimStatusDerivation\ A factorisation algebra on $\mathbb{C}$ with holomorphic
locality is equivalent to a vertex algebra (Costello--Gwilliam 2017 Vol I Thm 5.3.3).

\ClaimStatusConjectured\ $\mathcal{F}_{4D \mathcal{N} = 2}^{U(1), \Omega}$ as a vertex
algebra is isomorphic to $\Wilpha[\lambda]$ at $\lambda = \epsilon_1/\epsilon_3 +
\epsilon_2/\epsilon_3$ and central charge $c = 2 + (\epsilon_1 - \epsilon_2)^2/(\epsilon_1 \epsilon_2)$
or similar expression.
\emph{Status:} consistent with OPE computations at low order;
partial matchings in Costello 2013 §11 and Costello--Gaiotto 2018
\texttt{arXiv:1810.01970}. A full identification theorem is in the Gaiotto--Rapčák
program and Costello--Paquette 2020 \texttt{arXiv:2009.04834}.

\subsubsection{Holomorphic Chern--Simons on $\mathbb{C}^3$: the BV--BRST formulation}
\label{wn:subsubsec:hCS-on-C3}

\ClaimStatusDefinition\ \emph{hCS action.} Let $X = \mathbb{C}^3$ with holomorphic
volume form $\Omega = dz_1 \wedge dz_2 \wedge dz_3$. Take gauge group $G = U(1)$
(Abelian) for concreteness.
The field is $\mathcal{A} \in \Omega^{0, \bullet}(X, \mathfrak{g})$, the Dolbeault
super-field of ghosts + gauge potential + antifields, with total grading (ghost + form).

Action:
\[
S_\hCS(\mathcal{A}) = \int_X \Omega \wedge \Big( \mathcal{A} \bar\partial \mathcal{A}
+ \tfrac{2}{3} \mathcal{A} \wedge [\mathcal{A}, \mathcal{A}] \Big).
\]
For Abelian $G$, the cubic term vanishes and hCS is Gaussian.

\ClaimStatusTheorem\ (Costello--Gwilliam 2017 Vol II §5.1.)
The space of observables $\mathrm{Obs}(\hCS)$ is a factorisation algebra on $X = \mathbb{C}^3$
with respect to the Dolbeault-factorisation structure.

\ClaimStatusDefinition\ \emph{Line defect / spectral curve.}
Introduce a defect along a complex line $\mathbb{C}_z \subset \mathbb{C}^3$. The defect
supports its own observable algebra $\mathcal{F}_{\mathrm{defect}}$ which is a
factorisation algebra on $\mathbb{C}_z$, a module for the bulk $\mathrm{Obs}(\hCS)$ via
restriction.

\ClaimStatusConjectured\ (Costello 2013 §11 + Costello--Gaiotto 2018
\texttt{arXiv:1810.01970} framework.)
$\mathcal{F}_{\mathrm{defect}} \cong \Wilpha[\lambda]$ as factorisation algebras
(equivalently, vertex algebras) on the defect curve, with $\lambda$ determined by
the transverse equivariant weights.

\emph{What has been derived.} The BV--BRST complex of hCS has been computed explicitly
in Costello 2013 §9--10, reducing to free-boson / tower of differential
polynomials in the Abelian case. The identification with $\Wilpha$ at low orders
(central charge, first few generating fields) is verified; the full identification
requires the Gaiotto--Rapčák / Costello--Paquette framework.

\subsubsection{$E_3$-algebra structure of Chern--Simons observables}
\label{wn:subsubsec:E3-CS-2026}

\ClaimStatusTheorem\ (Costello--Gwilliam 2017 Vol I Thm 5.3.3 for the $E_n$-algebra
aspect.) 3D Chern--Simons theory on $\mathbb{R}^3$ has factorisation-algebra observables
forming an $E_3$-algebra (locally constant factorisation algebra = $E_n$-algebra on
$\mathbb{R}^n$ by Lurie's theorem \emph{Higher Algebra} §5.5.3).

\emph{Scope note.} The claims below rest on two verified sources: Costello--Gwilliam 2017
for the factorisation-algebra $E_n$-algebra framework, and Costello 2013
\texttt{arXiv:1303.2632} for the gauge-theory construction of the Yangian via
$\Omega$-background. A further-developed Costello--Francis--Gwilliam $E_3$-deformation
analysis of Chern--Simons theory sharpening the deformation-parameter story is, at the
time of this treatise, not in hand as a cited primary source; this treatise does not
rely on it.

\ClaimStatusDerivation\ \emph{Six-dimensional hCS on $\mathbb{C}^3 = \mathbb{R}^6$.}
hCS on $\mathbb{C}^3$ satisfies holomorphic locality: $\mathrm{Obs}(\hCS)|_{\mathrm{disk}}$
depends on the Dolbeault structure of the disk.
The observable algebra is neither $E_3$ (topological) nor $E_1$ (just associative),
but rather a \emph{holomorphic $E_3$-algebra}, i.e.~a factorisation algebra on $\mathbb{C}^3$
with holomorphic rather than topological locality.

\emph{Deformation parameter.} In the BV--BRST complex, $\hbar$ is the
BV/loop formal parameter, not a cohomological degree shift; the shifts
come from the BV and cochain conventions.  The $\Omega$-background gives
additional parameters $\epsilon_1,\epsilon_2,\epsilon_3$ appearing in the
transverse weight structure. Costello--Paquette 2020
\texttt{arXiv:2009.04834} applies the $E_n$-algebra BV formalism to the
4D $\mathcal{N}=2$ / vertex-algebra setup.  Costello--Francis--Gwilliam
2026 concerns ordinary three-dimensional Chern--Simons factorisation
algebras and knot polynomials; it is not a primary source for
six-dimensional hCS on a CY$_3$.  The hCS $E_3$ language here is an
avatar comparison built from Costello--Gwilliam factorisation algebras,
Costello's gauge-theory Yangian construction, and the holomorphic
BV--BRST complex.

%--------------------------------------------------------------------------------
\subsection{Example 2: Resolved conifold}
\label{wn:subsec:resolved-conifold}
%--------------------------------------------------------------------------------

\subsubsection{Geometry and quiver}

\ClaimStatusDefinition\ \emph{Resolved conifold.}
$\mathbb{Y} := \mathrm{Tot}(\mathcal{O}_{\mathbb{P}^1}(-1) \oplus \mathcal{O}_{\mathbb{P}^1}(-1)) \to \mathbb{P}^1$.
This is a local (non-compact) toric Calabi--Yau three-fold: the total space of two
line bundles over $\mathbb{P}^1$ with total Chern class $c_1 = -2$ matching
$c_1(K_{\mathbb{P}^1}) = -2$ by CY$_3$ adjunction.

\ClaimStatusDefinition\ \emph{Quiver $Q_{\mathrm{con}}$.}
Two vertices $\circ, \bullet$ (corresponding to the two toric fixed points on $\mathbb{P}^1$).
Four arrows: $a_1, a_2: \circ \to \bullet$ and $b_1, b_2: \bullet \to \circ$.

\ClaimStatusDefinition\ \emph{Potential $W = \mathrm{tr}(a_1 b_1 a_2 b_2 - a_1 b_2 a_2 b_1)$.}

\ClaimStatusTheorem\ (Szendroi 2008 \emph{Non-commutative Donaldson--Thomas theory of
the conifold}, Geom. Topol. 12.) The Jacobi algebra $J(Q_{\mathrm{con}}, W)$ is isomorphic
as an $R$-algebra to the non-commutative crepant resolution of the conifold singularity
(for $R$ the local ring of the singularity), and its derived category is equivalent to
$D^b(\mathbb{Y})$.

\subsubsection{The CoHA of the conifold}

\ClaimStatusDefinition\ Moduli of representations of $(Q_{\mathrm{con}}, W)$ with dimension
vector $(d_1, d_2)$: $\mathcal{M}_{d_1, d_2} = [\mathrm{Rep}(Q_{\mathrm{con}}, d_1, d_2) / \prod GL_{d_i}]$.

\ClaimStatusDefinition\ $\CoHA(\mathbb{Y}) := \bigoplus_{(d_1, d_2)} H^*_T(\mathcal{M}_{d_1, d_2}, \phi_W)$.

\ClaimStatusTheorem\ (Davison 2012 \texttt{arXiv:1209.4620} for conifold specifically;
Morrison--Mozgovoy--Nagao--Szendroi 2010 for DT-invariants of the conifold.)
\[
\CoHA(\mathbb{Y}) \cong U(\mathfrak{g}_{\mathrm{BPS}}(\mathbb{Y}))
\]
as associative algebras, where $\mathfrak{g}_{\mathrm{BPS}}(\mathbb{Y})$ is a specific
Lie algebra determined by the conifold geometry.

\subsubsection{Positive half generators: explicit form}

\ClaimStatusDerivation\ \emph{Dimension vector $(1, 0)$.}
A representation is a pair (single vector space on $\circ$, zero on $\bullet$, no arrows).
Moduli: $\mathcal{M}_{1, 0} = [\mathrm{pt}]$. Contribution: $\mathbb{F}$.

\emph{Dimension vector $(0, 1)$.} Similarly $\mathbb{F}$. Generator $f_0$.

\emph{Dimension vector $(1, 1)$.}
Two vector spaces, arrows $a_1, a_2 \in \mathbb{C}$ (from $\circ$ to $\bullet$) and
$b_1, b_2 \in \mathbb{C}$ (back). Modulo potential $W$: $a_1 b_1 a_2 b_2 - a_1 b_2 a_2 b_1$.
For $1 \times 1$ matrices this commutator vanishes identically, so the Jacobi
constraint is trivial. Moduli: $[\mathbb{A}^4 / (\mathbb{C}^\times)^2]$ (divide by each
$GL_1$).
Equivariant cohomology: $H^*_T(\mathbb{A}^4)^{T'}$ where $T' = (\mathbb{C}^\times)^2/\Delta$.

\emph{Generators at $(1, 1)$:} two ``hop'' generators, one for each pair of arrows,
giving $e_0, e_1$ (positive half) and $f_0, f_1$ (negative half) of the conifold Yangian.

\ClaimStatusTheorem\ (Rapčák--Soibelman--Yang--Zhao 2020 \texttt{arXiv:2001.10549} Thm~B:
cohomological conifold CoHA; Kapranov--Vasserot 2018 \texttt{arXiv:1802.07988} for the
K-theoretic quantum toroidal $\mathfrak{gl}_2$ incarnation.) The positive-half CoHA of the
conifold matches the positive half of the quantum toroidal algebra
$U_{q_1, q_2}(\widehat{\widehat{\mathfrak{gl}}_2})$ at a specific point, where
$q_1 = e^{\epsilon_1}, q_2 = e^{\epsilon_2}$ are the two equivariant parameters of
$T^2 \curvearrowright \mathbb{Y}$; the cohomological side is Rapčák--Soibelman--Yang--Zhao,
the K-theoretic toroidal identification is Kapranov--Vasserot.

\subsubsection{Shuffle formula for conifold CoHA}

\ClaimStatusDefinition\ \emph{Shuffle kernel for conifold.}
Over $\mathbb{F}[\epsilon_1, \epsilon_2]$ with the third weight on the compact $\mathbb{P}^1$
direction being $\epsilon_3 = 0$ (the CY$_3$ condition is $\epsilon_1 + \epsilon_2 = 0$
restricted to the fibre directions only; the $\mathbb{P}^1$-direction is compact and does
not enter the equivariant parameter):
\[
\omega_{\mathrm{con}}(z, w) = \frac{(z - w + \epsilon_1)(z - w + \epsilon_2)}{(z - w)(z - w + \epsilon_1 + \epsilon_2)}.
\]
(Negut 2015 \texttt{arXiv:1505.01528}.)

\ClaimStatusDerivation\ \emph{$(1, 1) \star (1, 0)$ multiplication.}
$(e_0 \star e_0')(z, w) = e_0(z) e_0'(w) \omega_{\mathrm{con}}(w, z) + e_0(w) e_0'(z) \omega_{\mathrm{con}}(z, w)$
(direct from shuffle formula).

\subsubsection{Chiral Yangian from conifold}

\ClaimStatusConjectured\ (following Costello--Li 2016 \texttt{arXiv:1606.00365},
Costello--Gwilliam Vol II §6 conifold example.)
The holomorphic Chern--Simons on $\mathbb{Y}$ gives a factorisation algebra
$\mathcal{F}_\hCS(\mathbb{Y})$ on $\mathbb{Y}$ (or on a chosen defect curve).

When restricted to a defect line $\mathbb{C}_z \subset \mathbb{Y}$ along a generic fibre
direction, the defect factorisation algebra $\mathcal{F}_{\mathrm{defect}, \mathrm{con}}$
is a vertex algebra. It is conjectured to be a quantum toroidal algebra at specific
parameters, or equivalently a deformation of $\Wilpha$ corresponding to the conifold
topology.

\ClaimStatusOpen\ \emph{Explicit vertex-algebra generators and OPEs of the conifold
chiral Yangian.} The programme of Costello--Li 2016 + Costello--Paquette 2020 provides
the setup but explicit formulas for all generators and OPEs are not, to my knowledge,
a single-citation result. Awata--Feigin--Shiraishi 2011
\texttt{arXiv:1112.6074} provide the relevant quantum toroidal structure.

%--------------------------------------------------------------------------------
\subsection{Example 3: $K3 \times E$}
\label{wn:subsec:K3xE}
%--------------------------------------------------------------------------------

\providecommand{\kappaBKM}{\kappa_{\mathrm{BKM}}}
\providecommand{\kappach}{\kappa_{\mathrm{ch}}}
\providecommand{\kappacat}{\kappa_{\mathrm{cat}}}
\providecommand{\Auts}{\mathrm{Aut}_s}
\providecommand{\Mred}{\mathcal{M}^{\mathrm{red}}_{DT}}
\providecommand{\Meff}{\mathcal{M}^+_{\mathrm{eff}}}

\subsubsection{Positive effective geometry via reduced Donaldson--Thomas}
\label{wn:subsubsec:K3xE-positive-geometry-reduced-DT}

\begin{theorem}[Positive effective geometry of $K3\times E$ via Pandharipande--Oberdieck reduced DT]
\label{wn:thm:K3xE-positive-geometry-reduced-DT}
Let $X = K3 \times E$ and let $\gamma \in H^{\mathrm{ev}}(X, \mathbb{Z})$ restrict to a
primitive curve class on the $K3$ factor. The positive effective moduli stack carrying
the CoHA of $X$ is the Pandharipande--Oberdieck reduced Donaldson--Thomas stack
\[
\Meff(K3 \times E) \;=\; \Mred(K3 \times E;\, \gamma\text{ primitive on }K3),
\]
equivariant for the rigid two-factor group
\[
G_{K3 \times E} \;=\; \mathbb{C}^\times_E \,\times\, \Auts(K3),
\]
with $\mathbb{C}^\times_E$ the translation $\mathbb{C}^\times$ on the elliptic factor and
$\Auts(K3) \subset M_{23}$ the symplectic-automorphism group of a Mukai lattice polarisation.
Its reduced DT character is
\[
Z^{\mathrm{red}}_{DT}(K3 \times E)(q, t, p)
\;=\; -\,\frac{1}{\Phi_{10}(q, t, p)}
\;=\; -\,\Delta_5^{-2}(q, t, p),
\]
where $\Phi_{10}$ is the Igusa cusp form (Oberdieck--Pixton 2017,
\texttt{arXiv:1706.10100}, \emph{Invent.\ Math.} 222 (2020))
and $\Delta_5$ is the Gritsenko lift of $\phi_{0,1}$
(Gritsenko--Nikulin 1998, \emph{Amer.\ J.\ Math.} 119). The associated Borcherds-Kac-Moody
weight is
\[
\kappaBKM(\Delta_5) \;=\; \frac{c_1(0)}{2} \;=\; 5.
\]
\end{theorem}

\begin{proof}[Proof sketch]
The reduced virtual class on $\mathcal{M}_{DT}(K3 \times E)$ exists by
Maulik--Pandharipande--Thomas 2010 (\texttt{arXiv:1001.2719}): the $E$-translation action
produces a holomorphic symplectic form whose dual trivialises one copy of the obstruction
bundle, yielding a class of virtual dimension one higher than the naive DT class. The
$\gamma$-primitive sector is $\mathbb{C}^\times_E$-equivariant via this translation, and
the symplectic-automorphism group $\Auts(K3)$ permutes the Mukai-lattice orbit, giving the
stated equivariance. The partition-function identity
$Z^{\mathrm{red}}_{DT} = -1/\Phi_{10}$ is the Oberdieck--Pixton theorem in the primitive
class. The Borcherds-product expansion of $\Phi_{10} = \Delta_5^2$
(Gritsenko--Nikulin 1998; Borcherds 1998 \emph{Invent.\ Math.} 132) gives the BKM weight
$\kappaBKM(\Delta_5) = c_1(0)/2 = 5$, the Fourier coefficient of $q^0$ in the weight-$5$
theta lift being $c_1(0) = 10$.
\end{proof}

\begin{remark}[Four-corner commutative diagram]
\label{wn:rem:K3xE-four-corner}
The reduced-DT positive geometry sits at the geometry corner of a commutative square
\[
\begin{array}{ccc}
\Mred(K3 \times E;\, \gamma\text{ primitive}) & \longleftrightarrow & \CoHA^+_{Y}(K3 \times E) \\
\big\updownarrow & & \big\updownarrow \\
\mathfrak{g}_{\mathrm{BPS}}(K3 \times E) & \longleftrightarrow & \mathfrak{g}_{\Delta_5}
\end{array}
\]
with the four corners: (i) the Pandharipande--Oberdieck geometry; (ii) the positive
half $Y^+(K3 \times E)$ of the would-be $K3 \times E$-Yangian, built as the CoHA on
$\Mred$; (iii) the BPS Lie algebra $\mathfrak{g}_{\mathrm{BPS}}(K3 \times E)$
(Davison 2017 \texttt{arXiv:1601.02479}; Davison--Meinhardt 2020
\emph{Invent.\ Math.} 221) of primitive elements; (iv) the Gritsenko--Nikulin BKM
superalgebra $\mathfrak{g}_{\Delta_5}$. The horizontal arrows are the CoHA-as-universal
enveloping algebra (Davison--Meinhardt PBW) and the Borcherds construction
($\mathfrak{g}_{\Delta_5}$ as the BKM with denominator $\Delta_5$). The vertical left
arrow is Davison's $\Omega$-invariant construction; the vertical right arrow is the
denominator-formula multiplicity identification. The square commutes at the level of
graded dimensions unconditionally; see
Remark~\ref{wn:rem:K3xE-open-problem-reconciliation} for the bracket-level status.
\end{remark}

\begin{remark}[Why toric localisation fails and what replaces it]
\label{wn:rem:K3xE-no-torus-replaced}
Generic $K3$ admits no non-trivial $\mathbb{C}^\times$-action: the classical
Bogomolov--Beauville decomposition forces a symplectic holomorphic form, and a
$\mathbb{C}^\times$-action would contract it. The $\mathbb{C}^3$-strategy of
Maulik--Okounkov / Schiffmann--Vasserot --- torus-fixed-point localisation on a toric
CY$_3$ --- therefore does not apply. The replacement comprises three distinct
geometric inputs:
\begin{enumerate}
\item \emph{Reduced virtual class} (Maulik--Pandharipande--Thomas 2010,
\texttt{arXiv:1001.2719}): the $E$-symplectic form trivialises one copy of the
obstruction bundle on $K3 \times E$, producing a class $[\Mred]^{\mathrm{vir}}$ of
virtual dimension one higher than the naive $[\mathcal{M}_{DT}]^{\mathrm{vir}}$. This
is what makes the integrand non-zero on the primitive locus.
\item \emph{Residual $\mathbb{C}^\times_E$} from elliptic translation: the factor $E$
carries its own $\mathbb{C}^\times$-action by translation (equivalently, by the
addition of $E$ to itself fibrewise), and this is a genuine scaling on the
$\mathbb{C}^3$-direction tangent to $E$ in the moduli.
\item \emph{Discrete $\Auts(K3) \subset M_{23}$ orbit} (Mukai 1988): the
symplectic-automorphism group acts on $H^*(K3, \mathbb{Z})$ by lattice isometries
preserving the Mukai polarisation; equivariant cohomology with respect to this
discrete group is a genuine refinement, even in the absence of continuous
$\mathbb{C}^\times$-action on the $K3$-factor.
\end{enumerate}
The combined group $G_{K3 \times E} = \mathbb{C}^\times_E \times \Auts(K3)$ is what
replaces $T^3$ in the Maulik--Okounkov / Schiffmann--Vasserot toric template.
\end{remark}

\subsubsection{Setup and obstructions}

$X = K3 \times E$: compact smooth Calabi--Yau three-fold, $\dim_\mathbb{C} X = 3$.
$h^{0, 0} = 1, h^{1, 0} = 1, h^{0, 1} = 1, h^{2, 0} = 1, h^{1, 1} = 21,\ldots$
(Hodge diamond of $K3$ times that of $E$).

\emph{Obstructions to the $\mathbb{C}^3$ strategy:}
\begin{enumerate}
\item $X$ is \emph{not toric}. No algebraic torus $T^3 \curvearrowright X$.
Generic $K3$ has no non-trivial $\mathbb{C}^\times$-action at all. So equivariant
cohomology is trivial (= ordinary cohomology) and localisation reduces to identities.
\item The elliptic curve factor $E$ contributes an extra $H^*(E) = \mathbb{C} \oplus
\mathbb{C}[1]$ to the Hodge structure: the CY$_3$ is of ``rigid'' type, in a specific
sense.
\item The CoHA construction via $T^3$-equivariant cohomology does not apply directly.
\end{enumerate}

\subsubsection{Davison's approach: perverse sheaves and critical cohomology}

\ClaimStatusConjectured\ (Davison 2017 \texttt{arXiv:1601.02479};
Davison--Meinhardt 2020 \emph{Invent. Math.} 221.) \emph{Scope hypothesis:}
$X$ a smooth proper Calabi--Yau three-fold admitting a global critical-locus /
potential presentation for its sheaf moduli (not automatic for every compact CY$_3$;
Davison's construction is formulated for quiver-with-potential models and their direct
geometric counterparts, with the global critical-chart hypothesis load-bearing).
Under this hypothesis, there is a cohomological Hall algebra
\[
\CoHA^{\mathrm{BPS}}(X) := H^*(\mathcal{M}(X), \phi_W \cdot \mathrm{IC}_{\mathcal{M}(X)}),
\]
an associative algebra satisfying:
\begin{enumerate}
\item Integration map $\pi_*: \CoHA^{\mathrm{BPS}}(X) \to $ motivic Hall algebra.
\item Existence of a BPS Lie subalgebra $\mathfrak{g}_{\mathrm{BPS}}(X) \subset \CoHA^{\mathrm{BPS}}(X)$
of primitive elements.
\item PBW theorem: $\CoHA^{\mathrm{BPS}}(X) = U(\mathfrak{g}_{\mathrm{BPS}}(X))$ as
associative algebras.
\end{enumerate}

\ClaimStatusDefinition\ $\mathfrak{g}_{\mathrm{BPS}}(K3 \times E)$ is thus \emph{defined}
via Davison's construction. Its properties are the subject of the programme.

\subsubsection{Numerical identification with $\mathfrak{g}_{\Delta_5}$}

\ClaimStatusTheorem\ (Oberdieck--Pixton 2017 \texttt{arXiv:1706.10100},
\emph{Invent. Math.} 222 (2020), for \emph{primitive} K3 class.)
\[
Z^{\mathrm{red}}_{DT}(K3 \times E; \gamma \text{ primitive on } K3)(q, t, p)
= - \frac{C}{\Phi_{10}(q, t, p)}
\]
where $\Phi_{10}$ is the Igusa cusp form and $C$ is an explicit constant.
\emph{Scope restriction:} the theorem as stated covers the primitive-class reduced DT
partition function; the imprimitive / multiple-cover contribution is not inside the
stated identity.

\ClaimStatusDerivation\ \emph{Kontsevich--Soibelman wall-crossing formula.}
The reduced DT partition function admits an infinite-product expansion
\[
Z^{\mathrm{red}}_{DT} = \prod_{\gamma \in \Gamma^{\mathrm{eff}}}
(1 - x^\gamma)^{- (-1)^{|\gamma|}\Omega(\gamma)}
\]
where $\Omega(\gamma) \in \mathbb{Z}$ are the numerical BPS invariants.

\ClaimStatusTheorem\ (Gritsenko--Nikulin 1998 \emph{Amer. J. Math.} 119.)
$\Delta_5$ admits a Borcherds-product expansion:
\[
\Delta_5(Z) = e^{-2\pi i \langle \rho, Z\rangle}
\prod_{\alpha > 0}(1 - e^{-2\pi i \langle \alpha, Z\rangle})^{(-1)^{|\alpha|}\mathrm{mult}_{\mathrm{BKM}}(\alpha)}.
\]

\ClaimStatusDerivation\ \emph{Matching.} Equating $\Phi_{10} = \Delta_5^2$ with
$Z^{\mathrm{red}}_{DT}{}^{-1}$ (up to scalar $C$), and comparing infinite-product
coefficients charge-by-charge:
\[
\Omega(\gamma) = \mathrm{mult}_{\mathrm{BKM}}(\alpha_\gamma) \quad \text{for all } \gamma.
\]

\ClaimStatusDerivation\ \emph{Dimension equality.}
By Davison's construction $\Omega(\gamma) = \dim \mathfrak{g}_{\mathrm{BPS}, \gamma}$.
Therefore
\[
\dim \mathfrak{g}_{\mathrm{BPS}, \gamma}(K3 \times E) = \mathrm{mult}_{\mathrm{BKM}}(\alpha_\gamma) = \dim \mathfrak{g}_{\Delta_5, \alpha_\gamma}.
\]

\ClaimStatusOpen\ \emph{Lie-algebra isomorphism.}
Two Lie algebras with equal dimensions of weight spaces are not automatically isomorphic;
one must check that the brackets coincide. That is, for $\gamma_1 + \gamma_2 = \gamma$:
\[
[\mathfrak{g}_{\mathrm{BPS}, \gamma_1}, \mathfrak{g}_{\mathrm{BPS}, \gamma_2}] \to
\mathfrak{g}_{\mathrm{BPS}, \gamma} \quad \text{vs} \quad
[\mathfrak{g}_{\Delta_5, \alpha_{\gamma_1}}, \mathfrak{g}_{\Delta_5, \alpha_{\gamma_2}}]
\to \mathfrak{g}_{\Delta_5, \alpha_\gamma}.
\]
The first is determined by the CoHA product. The second is determined by the BKM
construction. Their equality — the Lie-algebra-level identification
$\mathfrak{g}_{\mathrm{BPS}}(K3 \times E) \cong \mathfrak{g}_{\Delta_5}$ — is an \emph{open}
problem. It is consistent with KS motivic wall-crossing but has not been reduced to
a single citation.

\subsubsection{Can we construct the positive-half CoHA and find the pairing for Drinfeld double?}

\ClaimStatusConjectured\ \emph{Positive half.} By analogy with $\mathbb{C}^3$ and the
conifold, one expects the positive half of a ``$K3 \times E$-Yangian'' to be realised on
some equivariant cohomology of a moduli space of sheaves on $K3 \times E$.

\emph{Obstacle:} the $T$-action is trivial on generic $K3$. However:
\begin{itemize}
\item Nikulin 1980 (1979 original) classifies $K3$ surfaces with finite
symplectic-automorphism groups. A \emph{Nikulin involution} $\iota: K3 \to K3$ fixing
the holomorphic 2-form gives a $\mathbb{Z}/2$-equivariant cohomology.
\item Generalising to Mukai's 11-group classification (Mukai 1988): all finite
symplectic automorphism groups of $K3$ are subgroups of $M_{23} \subset M_{24}$.
\item For each Mukai group $G$, $G$-equivariant cohomology $H^*_G(K3)$ is defined.
\item Combined with $E$-translation subgroup $E[N] \cong (\mathbb{Z}/N)^2$, one gets
$G \times E[N]$-equivariant cohomology of $K3 \times E$.
\end{itemize}

\ClaimStatusConjectured\ \emph{$G \times E[N]$-equivariant CoHA of $K3 \times E$.}
For programme-core $N \in \{1, 2, 3, 4, 6\}$, construct
\[
\CoHA^{G_N}(K3 \times E) := \bigoplus_\gamma H^*_{G_N \times E[N]}(\mathcal{M}(K3 \times E, \gamma), \phi_W)
\]
with multiplication via correspondences.

\emph{Status of construction:} this has not been carried out in the literature to my
knowledge. Davison 2017's framework is the closest theoretical basis, but explicit
computations on $K3 \times E$ await.

\ClaimStatusConjectured\ \emph{Pairing and Drinfeld double.}
Assuming the positive half exists, a bialgebra structure requires a non-degenerate
Hopf pairing. For the compact $K3 \times E$ setting, the pairing would come from
Serre duality on $D^b(K3 \times E)$:
\[
\langle \cdot, \cdot \rangle: \CoHA^{G_N}(K3 \times E)_+ \otimes \CoHA^{G_N}(K3 \times E)_- \to \mathbb{F}_{G_N}
\]
via $\mathrm{Ext}^3$-pairing on the moduli stacks. This is the ``Serre self-duality''
of the CoHA that Davison--Meinhardt 2020 use for their PBW theorem.

Given the pairing, Drinfeld doubling yields a Hopf algebra whose Cartan part
$Y^0_{K3 \times E}$ would be the ``programme-compatible'' completion.

\ClaimStatusOpen\ \emph{Does the Cartan part encode $\Delta_5$?}
The fundamental conjecture is: $Y^0_{K3 \times E}$ is generated by Cartan oscillators
$\psi_k$ whose generating series $\psi(z) = \exp(\sum_k h_k z^{-k-1})$ is a
$\Delta_5$-automorphic function in the sense that
\[
\psi(z)\big|_{z \mapsto Z\text{-coordinate on } \mathcal{A}_2} \propto \Delta_5^{-1}(Z)\text{-expansion.}
\]
This is the target of the programme but is not established.

\subsubsection{The positive half, explicit attempt}

Let us attempt the construction in programme-core case $N = 1$ (trivial $G_N$):

\ClaimStatusDefinition\ Moduli $\mathcal{M}_{K3 \times E}(\gamma)$ for $\gamma \in H^{\mathrm{ev}}(K3 \times E, \mathbb{Z})$
effective = $(r, \beta, n)$ with $r, n \in \mathbb{Z}$, $\beta \in H^2(K3 \times E, \mathbb{Z})$.

\ClaimStatusDefinition\ $\CoHA^+_{K3 \times E} := \bigoplus_\gamma H^*(\mathcal{M}(\gamma), \phi_W)$
(non-equivariant; or over $\mathbb{Q}$-algebraic cycles).

\ClaimStatusDerivation\ \emph{Associativity.} The KS correspondence structure gives
associativity. (Derived from general Hall-algebra principles.)

\ClaimStatusConjectured\ \emph{Generators.}
Low-rank low-charge generators are concentrated at:
\begin{itemize}
\item $\gamma = (1, 0, 0)$: ideal sheaf of a point on $K3 \times E$; moduli $= K3 \times E$.
Cohomology $H^*(K3 \times E) = H^*(K3) \otimes H^*(E)$, dimension $24 \cdot 4 = 96$.
\item $\gamma = (0, \beta, n)$: Pandharipande--Thomas stable pair with class $\beta$,
$n$ free points. At primitive $\beta$ on $K3$: moduli is a K3-fibred Hilbert scheme;
cohomology computed by Oberdieck--Pixton 2017.
\item $\gamma = (-1, 0, n)$: ``anti-D0-D6'' bound state; moduli comprises $n$-tuples of
points on $K3 \times E$ modulo the virtual structure.
\end{itemize}

\ClaimStatusDerivation\ \emph{Generating function.} The Hilbert series (Poincaré polynomial)
of $\CoHA^+_{K3 \times E}$, weighted by charges, should reproduce the ``positive-chamber''
contribution to $Z^{\mathrm{red}}_{DT}$. Explicitly, the character
\[
\chi(\CoHA^+_{K3 \times E})(q, t, p) = \sum_\gamma \dim H^*(\mathcal{M}(\gamma), \phi_W) \cdot x^\gamma
\]
should satisfy $\chi = Z^{\mathrm{red}}_{DT}|_{\mathrm{primitive}}$.

\emph{Status.} Order-by-order verified at low charges (Pandharipande--Thomas 2014 +
Oberdieck--Pixton 2017). Closed-form equivalence to $\Phi_{10}^{-1}$ on the primitive
K3 class is the Oberdieck--Pixton 2017 theorem (\emph{Invent. Math.} 222 (2020)).

\subsubsection{Localisation strategy}

The fact that generic $K3$ has no torus action does not mean localisation is
unavailable. Three alternative localisation strategies:

\ClaimStatusTheorem\ \emph{Strategy 1: Virtual localisation on $E[N]$.}
The translation subgroup $E[N] \cong (\mathbb{Z}/N)^2$ acts freely on $E$. The
quotient $X^{(N)} = (K3 \times E)/E[N]$ is a CHL orbifold (Chaudhuri--Hockney--Lykken
1995). Virtual-class localisation on the fixed locus of $E[N]$:
\[
[\mathcal{M}(X^{(N)})]^{\mathrm{vir}} = \frac{[\mathcal{M}(X^{(N)})^{E[N]}]^{\mathrm{vir}}}{e(N_{\mathrm{vir}})}.
\]

\ClaimStatusTheorem\ \emph{Strategy 2: Reduction to $K3$ fibres.}
Since $X = K3 \times E$, the projection $X \to E$ has K3 fibres. Pandharipande--Thomas 2014
compute the K3-fibre contribution explicitly; the $E$-modular dependence then
factors out.

\ClaimStatusTheorem\ \emph{Strategy 3: Nakajima construction.}
$H^*(\Hilb^n(K3))$ with its Heisenberg action gives a Fock-space realisation without
needing a torus action on $K3$; only the cohomology lattice of $K3$ is required.
This gives
\[
\bigoplus_n H^*(\Hilb^n(K3 \times E)) \cong \bigoplus_n H^*(\Hilb^n(K3)) \otimes H^*(E^{\otimes n}_{\text{Sym}})
\]
via the Göttsche formula extended to products.

\ClaimStatusOpen\ \emph{Integration of the three strategies into a single CoHA with
Drinfeld-double structure producing $\mathfrak{g}_{\Delta_5}$.} This is exactly the
unresolved programme question.

\subsubsection{What would a chiral BKM on $K3 \times E$ look like?}

By analogy with Costello 2013 for $\mathbb{C}^3$:

\ClaimStatusConjectured\ 4D $\mathcal{N} = 2$ theory obtained from M5-branes on
$K3 \times E$ (the Vafa--Witten / $(2, 0)$-theory setup), with $\Omega$-deformation
on $\mathbb{C}^2_{\epsilon_1, \epsilon_2}$ transverse to $K3 \times E$, compactified
along a defect curve $\Sigma \subset K3 \times E$ (e.g. the $E$-factor, or a hyperelliptic
curve in $K3$), gives a factorisation algebra on $\Sigma$.

\emph{Candidate $\Sigma$ choices:}
\begin{itemize}
\item $\Sigma = E$: the factorisation algebra on the elliptic curve is a ``chiral
algebra twisted by $E$-monodromy'', related to elliptic genera.
\item $\Sigma = \mathbb{C}$ after holomorphic-topological twist: a factorisation
algebra on the spectral line, analogous to Costello 2013.
\item $\Sigma = $ genus-2 hyperelliptic in $K3$: the Gaiotto curve at $N = 1$
(Wave-7 B1 identification, class-$\mathcal{S}$ $(A_1, \Sigma_{2, 0})$).
\end{itemize}

\ClaimStatusOpen\ \emph{Identification with $\mathfrak{g}_{\Delta_5}$-chiral-Yangian.}
The CY-to-chiral functor $\Phi_3: D^b(K3 \times E) \to E_1\text{-chiral}$ of the Vol III
programme, applied at $(K3 \times E, \mathrm{point})$, is conjectured to yield a chiral
algebra whose BPS / primitive sub-object is $\mathfrak{g}_{\Delta_5}$. This is the
programme's central conjecture, not a theorem.

\begin{remark}[Reconciliation of the ``open problem'' language]
\label{wn:rem:K3xE-open-problem-reconciliation}
The \texttt{ClaimStatusOpen} markers above
($\S$\ref{wn:subsec:K3xE} ``Lie-algebra isomorphism'', ``Does the Cartan part encode
$\Delta_5$?'', ``Integration of the three strategies'', ``Identification with
$\mathfrak{g}_{\Delta_5}$-chiral-Yangian'') refer precisely to the \emph{bracket-level}
identification
\[
\mathfrak{g}_{\mathrm{BPS}}(K3 \times E) \;\cong\; \mathfrak{g}_{\Delta_5}
\]
as Borcherds-Kac-Moody Lie superalgebras. This bracket identification remains
conditional on the Costello topological-conformal-field-theory cyclic-invariance input
(the $[m_k, B^{(2)}] = 0$ total-cocycle computation via TCFT/Gaiotto--Moore--Witten
pairing, not per-$k$), recorded as AP-CY34 in
\texttt{notes/antipatterns\_catalogue.md}.
By contrast, the \emph{graded-dimension} matching
\[
\dim \mathfrak{g}_{\mathrm{BPS}, \gamma}(K3 \times E)
\;=\; \dim \mathfrak{g}_{\Delta_5, \alpha_\gamma}
\qquad \forall\,\gamma\text{ primitive on }K3
\]
is unconditional: it follows from Oberdieck--Pixton 2017 (\emph{Invent.\ Math.} 222
(2020)) for the reduced DT partition function
$Z^{\mathrm{red}}_{DT} = -1/\Phi_{10}$, combined with Davison--Meinhardt 2020
(\emph{Invent.\ Math.} 221) for $\Omega(\gamma) = \dim \mathfrak{g}_{\mathrm{BPS},\gamma}$
and the Borcherds denominator-formula expansion of $\Phi_{10} = \Delta_5^2$
(Borcherds 1998 \emph{Invent.\ Math.} 132; Gritsenko--Nikulin 1998
\emph{Amer.\ J.\ Math.} 119) for $\dim \mathfrak{g}_{\Delta_5, \alpha_\gamma}$.
The positive effective geometry $\Meff(K3 \times E) = \Mred$ of
Theorem~\ref{wn:thm:K3xE-positive-geometry-reduced-DT} is the rigid locus on which the
graded-dimension theorem operates; the bracket-level question is what the remaining
``open'' markers track.
\end{remark}

%--------------------------------------------------------------------------------
\subsection{Summary}
\label{wn:subsec:coha-W-summary}
%--------------------------------------------------------------------------------

\begin{center}
\begin{tabular}{l|l|l|l}
\textbf{Step} & \textbf{$\mathbb{C}^3$} & \textbf{Conifold} & \textbf{$K3 \times E$} \\\hline
Quiver + potential & Jordan triple, $\mathrm{tr}(X[Y,Z])$ & 2-vertex 4-arrow, $W_{\mathrm{con}}$ & N/A (compact) \\
Jacobi $\cong \mathcal{O}_{\mathrm{CY3}}$ & Theorem & Theorem (Szendroi) & N/A \\
CoHA as associative & Theorem (KS) & Theorem (Davison 2012) & Theorem (Davison 2017) \\
Shuffle formula & Explicit (SV 2013) & Explicit (Negut 2015) & \textsc{open} \\
Positive half $Y^+$ & Shuffle-Hopf & Conifold toroidal & \textsc{conjectural} \\
Drinfeld double & Theorem (Tsymbaliuk 2017 Thm~1.1) & Partial (RSYZ 2020; KV 2018) & \textsc{open} \\
Identified with VOA & $\Wilpha$ [Costello] & Quantum toroidal [CL 2016] & \textsc{open} \\
Gauge theory origin & 4D $\mathcal{N}=2$ $U(1)$ $\Omega$-bg & idem (non-compact) & Conjectural class-$\mathcal{S}$ \\
hCS factorisation alg & Explicit (CG 2017) & Partial (CL 2016) & \textsc{open} \\
CFG 2026 $E_3$-CS source & 3D CS only; hCS avatar & 3D CS only; hCS avatar & \textsc{not a 6D hCS source} \\
BPS Lie algebra & $\widehat{\mathfrak{gl}}_1$ & Toroidal $\widehat{\widehat{\mathfrak{gl}}_2}$ & $\mathfrak{g}_{\Delta_5}$? \\
\end{tabular}
\end{center}

The programme's centre of gravity is the last column, last row: the conjectural
identification of the BPS Lie algebra on $K3 \times E$ with the Gritsenko--Nikulin
BKM $\mathfrak{g}_{\Delta_5}$, and the associated chiral Yangian.
The first two examples provide the concrete toolkit and the vocabulary; the third
remains a programme target, with multiple open structural steps enumerated above.

\subsection{Holomorphic Chern--Simons and categorical Hochschild: scoped comparison}
\label{subsec:hcs-vs-cat-hochschild}

\begin{definition}[hCS datum]
\label{def:hcs-datum-treatise}
An \emph{hCS datum} on a complex manifold $X$ of complex dimension $3$ is a triple $(X, \mathfrak{g}, \Omega_X)$ with $\mathfrak{g}$ a finite-type $L_\infty$-algebra carrying an invariant non-degenerate pairing $\langle-,-\rangle$ and $\Omega_X \in H^{3,0}(X)$ a nowhere-vanishing section. The classical BV field is $\mathcal{A} \in \Omega^{0,\bullet}(X, \mathfrak{g})[1]$; the action is
\[
S_{\mathrm{cl}}[\mathcal{A}] \;=\; \int_X \Omega_X \wedge \Big( \tfrac{1}{2} \langle \mathcal{A}, \bar\partial \mathcal{A} \rangle \;+\; \tfrac{1}{6} \langle \mathcal{A}, [\mathcal{A}, \mathcal{A}] \rangle \Big),
\]
with the BV master equation $\{S, S\} = 0$ equivalent to the $L_\infty$-Jacobi identity on $\mathfrak{g}$.
\end{definition}

\begin{definition}[CY$_3$ categorical datum]
\label{def:cy3-cat-datum-treatise}
A \emph{CY$_3$ categorical datum} is a pair $(\mathcal{C}, \eta)$ with $\mathcal{C}$ a smooth proper $\mathbb{Z}$-graded dg-category and $\eta \colon \HH_{-3}(\mathcal{C}) \to k$ a non-degenerate pairing inducing the Serre functor $S_{\mathcal{C}} = [3]$. The Hochschild cochain complex $\HH^\bullet(\mathcal{C})$ carries a Gerstenhaber bracket of degree $1 - 3 = -2$ and a graded-commutative cup product; the Serre pairing raises this to a shifted-symplectic $L_\infty$-structure.
\end{definition}

\begin{theorem}[Formal-locus $E_3$ structure at $d = 3$]\ClaimStatusCorrected
\label{thm:kt-formality-d3-treatise}
Assume that $\HH^\bullet(\mathcal{C})$ is formal as an $E_3$-algebra.
Then the underlying $L_\infty$-structure on $\HH^\bullet(\mathcal{C})$
lifts to an $E_3$-algebra on the formal locus.  Operad-level
$E_3$-formality is unconditional over $\mathbb Q$; category-level
formality of this Hochschild algebra is the additional hypothesis.  One
explicit local model is the Kontsevich graph formula
\[
m_n(\phi_1, \ldots, \phi_n) \;=\; \sum_{\Gamma \in \mathcal{G}_{n, \mathrm{adm}}} w_\Gamma \cdot B_\Gamma(\phi_1, \ldots, \phi_n), \qquad w_\Gamma = \int_{\overline{\mathrm{Conf}}_n(\mathbb{R}^3)} \omega_\Gamma,
\]
with $\overline{\mathrm{Conf}}_n(\mathbb{R}^3)$ the
Fulton--MacPherson compactification and $\omega_\Gamma$ the wedge of
angle forms associated to edges.  Globalisation from formal disks to a
compact $X$ picks up the Calaque--Van den Bergh Duflo square-root
$\mathrm{td}(X)^{1/2}$ and the Atiyah-class obstruction.  Thus
$\mathbb{C}^3$ and the formal $K3\times E$ locus are in scope here; a
general compact CY$_3$ is not asserted to be formal.
\end{theorem}

\begin{theorem}[HKR on $\mathbb{C}^3$ identifies the two grammars]
\label{thm:hkr-c3-identifies-treatise}
Let $\mathcal{C} = D^b(\Coh(\mathbb{C}^3))$. The Caldararu--Willerton-twisted HKR map
\[
\mathrm{HKR}_{\mathrm{tw}} \colon \HH^\bullet(\mathcal{C}) \;\xrightarrow{\simeq}\; \PV^\bullet(\mathbb{C}^3) \;=\; \Omega^{0,\bullet}(\mathbb{C}^3, \wedge^\bullet T_{\mathbb{C}^3})
\]
is an isomorphism of $E_3$-algebras. For any compact generator $V \in \mathcal{C}$ of rank $r$, the derived endomorphism $\mathrm{dg}$-algebra is
\[
R\End^\bullet_{\mathcal{C}}(V) \;=\; \mathfrak{gl}_r \otimes \PV^\bullet(\mathbb{C}^3)
\]
as $A_\infty$-algebras, and the classical observables of hCS on $\mathbb{C}^3$ with gauge $\mathfrak{gl}_r$ are
\[
\mathrm{Obs}^{\mathrm{cl}}_{\hCS, \mathfrak{gl}_r}(\mathbb{C}^3) \;=\; \Omega^{0,\bullet}(\mathbb{C}^3, \mathfrak{gl}_r)[1],
\]
which agrees with $\mathfrak{gl}_r \otimes \PV^\bullet(\mathbb{C}^3)$ as $E_3$-algebras. The Bochner--Martinelli propagator
\[
P_{\mathrm{BM}}(z, w) \;=\; \frac{2}{(2\pi i)^3} \sum_{k=1}^{3} \frac{(-1)^{k-1} \overline{(z_k - w_k)}}{\|z - w\|^6} \, \widehat{d\bar z_k} \wedge dw_1 \wedge dw_2 \wedge dw_3
\]
gives the OPE $\mathcal{A}(z) \mathcal{A}(w) \sim \hbar P_{\mathrm{BM}}(z, w) \cdot \mathbf{1}$ modulo $Q$-exact; the Costello--Li heat-kernel regularisation degenerates to $P_{\mathrm{BM}}$ at $\epsilon \to 0$, and the resulting Feynman weights equal the Kontsevich configuration-space integrals, picking the Kontsevich associator in the $\mathrm{GRT}_1$-torsor.
\end{theorem}

\begin{theorem}[Three cocycles from the Atiyah class on compact $X$]
\label{thm:atiyah-three-cocycles-treatise}
For $X$ compact CY$_3$, the Atiyah class $\mathrm{At}(T_X) \in H^1(X, \Omega^1_X \otimes \End(T_X))$ sources three distinct cocycles which coexist on compact non-flat $X$ and which enter hCS and the categorical $E_3$-algebra at different orders:
\begin{enumerate}
\item \emph{Generator.} $\mathrm{At}(T_X) \in H^1(X, \Omega^1_X \otimes \End T_X)$: the $L_\infty$ binary bracket $\ell_2$ of the Kapranov $L_\infty$-algebra on $T_X[-1]$ (Kapranov 1999). Under the Calaque--C\u ald\u araru--Tu cochain-level HKR, $\mathrm{At}(T_X)$ is the cohomological defect from HKR being a Gerstenhaber morphism.
\item \emph{Tree-level Yukawa.} $Y_3 \in H^{0,3}(X) \simeq \mathrm{Sym}^3 H^1(X, T_X)^\vee$: the triple cup product of Atiyah classes against $\Omega_X$,
\[
Y_3(v_1, v_2, v_3) \;=\; \int_X \Omega_X \wedge \mathrm{At}(v_1) \cup \mathrm{At}(v_2) \cup \mathrm{At}(v_3),
\]
the Kapranov $\ell_3^{\min}$-bracket of the $L_\infty$-structure on $\wedge^\bullet T_X$. Physically: the Bershadsky--Cecotti--Ooguri--Vafa tree-level three-point Yukawa coupling.
\item \emph{One-loop BCOV curving.} $\alpha_{\mathrm{BCOV}} \in H^{0,1}(X)$: the one-loop BV anomaly
\[
\alpha_{\mathrm{BCOV}} \;=\; \frac{\chi(X)}{24} \cdot [\Omega_X]^{0,1}
\]
(Costello--Li 2016, Prop.~5.2). Under the CY isomorphism $H^{0,1}(X) \simeq H^{0,3}(X)^\vee$, $\alpha_{\mathrm{BCOV}}$ is Serre-dual to $Y_3$ but inhabits a distinct Hodge group.
\end{enumerate}
The $E_3$-curving of hCS observables is $\{Y_3, -\}$ entering the BV action as the cubic correction $\int_X \Omega_X \wedge Y_3 \wedge \langle \mathcal{A}, \mathcal{A}, \mathcal{A} \rangle$ at tree level; the $\alpha_{\mathrm{BCOV}}$-term enters as a one-loop BV anomaly independent of the tree-level cubic. The Atiyah class is the common source; $Y_3$ and $\alpha_{\mathrm{BCOV}}$ are distinct cocycles.
\end{theorem}

\begin{remark}[Status at $X = K3 \times E$]
\label{rem:three-cocycles-K3xE-treatise}
At $X = K3 \times E$:
\[
 \dim H^{0,3}(K3 \times E) \;=\; \dim\!\big(H^{0,2}(K3) \otimes H^{0,1}(E)\big) \;=\; 1 \cdot 1 \;=\; 1,
 \qquad
 \dim H^{0,1}(K3 \times E) \;=\; 1.
\]
The Hodge receptacles for $Y_3$ and the BCOV one-loop class are both
one-dimensional, but the formal-locus statement for $K3\times E$ is a
vanishing statement about the obstruction classes, not about the ambient
Hodge groups.  The elliptic factor has trivial tangent Atiyah class, the
K3 factor has no $\Omega^3_{K3}$ receptacle for the Kuranishi cubic, and
$\chi(K3\times E)=24\cdot0=0$ kills the BCOV coefficient.  Thus this
comparison records possible cocycle slots, while the $K3\times E$
Stage-$1$ object used in the main programme lies on the formal locus.
\end{remark}

\begin{theorem}[One-loop anomaly $\kappa_{\mathrm{anom}}$]\ClaimStatusCorrected
\label{thm:one-loop-anomaly-treatise}
The old cubic-Casimir anomaly formula is retracted.  The one-loop BV
analysis splits the logarithmic wave-function counterterm from the
cohomological gauge anomaly:
\[
 A_{\mathrm{w.f.}}(\mathfrak{g}) \;=\; -\frac{C_2(\mathfrak{g})}{(2\pi)^3} \;=\; -\frac{2 h^\vee}{(2\pi)^3}
\]
The quadratic-Casimir coefficient $A_{\mathrm{w.f.}}$ is the
wave-function renormalisation coefficient of the heat-kernel regulator
and is absorbed into a BV-trivial counterterm; it carries no
cohomological obstruction.  The cohomological holomorphic gauge anomaly
in complex dimension $d$ is instead governed by an invariant polynomial
of degree $d+1$ in the gauge curvature.  In the present $d=3$ case the
anomaly is quartic, schematically
\[
\kappa_{\mathrm{anom}}(X,\mathfrak g)
\;\sim\;
\hbar\int_X \Omega_X\wedge I_4(\mathcal A,F_\mathcal A,F_\mathcal A,F_\mathcal A),
\]
where $I_4\in\mathrm{Sym}^4(\mathfrak g^\vee)^{\mathfrak g}$ is the
quartic Chern--Weil colour factor in the chosen representation.  Three
consequences replace the cubic list:
\begin{enumerate}
\item \emph{Cubic vanishing is not enough.} Vanishing of
$d^{abc}$ for $\mathfrak{u}(1)$, $\mathfrak{sl}_2$, $\mathfrak{so}_N$,
or exceptional types does not by itself cancel the $d=3$ hCS anomaly;
one must evaluate the quartic invariant $I_4$.
\item \emph{$K3\times E$ kills the BCOV coefficient.}
$\chi_{\mathrm{top}}(K3\times E)=24\cdot0=0$ by Kunneth, so the
Euler-characteristic BCOV coefficient vanishes.  This is a base-space
vanishing, not a Lie-theoretic cubic-Casimir theorem.
\item \emph{The quintic remains a test case, not the old number.}
$\chi_{\mathrm{top}}(X_5)=-200$ makes the quintic the first compact
CY$_3$ test for the quartic anomaly, but the old coefficient
$-400N\hbar/(2\pi)^3$ was a cubic-Casimir artefact and is not retained.
\end{enumerate}
\end{theorem}

\begin{theorem}[DT and GW as conditional $E_3$-trace comparison]\ClaimStatusConjectured
\label{thm:DT-GW-MNOP-treatise}
On the formal-locus surface where $\mathcal{F}_X=\PhiFA_3(\mathcal{C})$
has been constructed, the Donaldson--Thomas and Gromov--Witten
partition functions may be interpreted as two trace functionals on the
same candidate factorisation algebra:
\begin{align*}
Z^{\mathrm{DT}}(q) \;&=\; \mathrm{Tr}^{\mathcal{F}_X}_{\mathrm{cat}} (\mathrm{id}) \;=\; \sum_{\gamma \in K_0(\mathcal{C})} \Omega_{\mathrm{BPS}}(\gamma) \cdot q^{\chi(\gamma)}, \\
Z^{\mathrm{GW}}(u) \;&=\; \mathrm{Tr}^{\mathcal{F}_X}_{\mathrm{geom}} (\mathrm{id}) \;=\; \sum_{g \geq 0} u^{2g - 2} F_g(t).
\end{align*}
The Maulik--Nekrasov--Okounkov--Pandharipande correspondence
\[
Z^{\mathrm{DT}}(q) \big|_{-q = e^{i u}} \;=\; C(u) \cdot Z^{\mathrm{GW}}(u)
\]
asserts that the two enumerative theories agree modulo the MacMahon
factor $C(u)=M(-q)^{\chi_{\mathrm{top}}(X)}$.  The trace language is a
conditional reinterpretation, not a proof of MNOP and not by itself an
hCS-versus-Hochschild comparison theorem.  At $X=K3\times E$ the DT side
reads $Z^{\mathrm{DT}}=C'\cdot\Delta_5(2Z)^{-2}$
(Oberdieck--Pandharipande 2018); the GW side follows by the
$-q=e^{iu}$ substitution.
\end{theorem}

\begin{corollary}[hCS and categorical Hochschild comparison, scoped]\ClaimStatusConjectured
\label{cor:hcs-geometric-resolution-treatise}
The categorical grammar $(\mathcal{C},\eta)$ supplies a candidate
$E_3$-structure on $\HH^\bullet(\mathcal{C})$ only under the
formal-locus and framing hypotheses above.  The geometric grammar
$(X,\mathfrak g,\Omega_X)$ supplies a computational model for a chosen
framing: the Bochner--Martinelli propagator, the Feynman graph expansion,
wave-function counterterms, and the corrected quartic anomaly slot.  The
two grammars differ in input type and are not literally two
presentations of one object without a specified HKR/KT quasi-isomorphism,
compatible framing, and factorisation-algebra lift.  Where those data
are supplied, the comparison is a construction target for
$\PhiFA_3(\mathcal C)$.
\end{corollary}

\begin{remark}[Per-grammar invariants]
\label{rem:per-grammar-invariants-treatise}
The asymmetry in visible invariants is Yoneda-asymmetry between an object and its framings: an external gauge $\mathfrak{g}$ is the endomorphism algebra $H^0(R\End_{\mathcal{C}}(V))$ of a chosen framing $V \in \mathcal{C}$; the whole $R\End_{\mathcal{C}}(-)$ sheaf over $\mathcal{C}^{\cong}$ is the internal datum. The geometric grammar fixes one fibre of this sheaf; the categorical grammar carries all fibres simultaneously, together with the 2-groupoid $\mathrm{Aut}^\otimes(\mathcal{C})$ of Fourier--Mukai autoequivalences that relate them. Morita-invariant categorical invariants (Serre pairing, Hochschild cohomology, formal-locus $E_3$-structure) are common to all framings; framing-specific invariants (quartic anomaly slot, BV determinant, specific propagator regularisation) depend on the choice of fibre.
\end{remark}

%--------------------------------------------------------------------------------
\subsection{The $E_3$ holomorphic algebra of observables: first-principles avatar}
\label{subsec:E3-algebra-of-observables-avatar}
%--------------------------------------------------------------------------------

Costello--Francis--Gwilliam 2026 (CFG) treats the ordinary three-dimensional Chern--Simons theory on $\R^3$: its factorisation algebra of observables is a locally-constant factorisation algebra on $\R^3$, hence by Lurie's recognition theorem an $E_3$-algebra in chain complexes. The present subsection constructs the six-dimensional holomorphic Chern--Simons avatar side by side with that source: the field-theoretic input is different (holomorphic rather than topological), but the operadic scaffolding is the same.

\subsubsection{The holomorphic $E_3$ operad on $\C^3$}
\label{sssec:holE3-operad}

\begin{definition}[Holomorphic $E_3$ operad]
\label{def:holE3-operad-avatar}
Let $\C^3$ be equipped with its Dolbeault complex structure. The operad $E_3^{\hol}$ of \emph{holomorphic little $3$-polydisks} has:
\begin{itemize}
\item $\operatorname{Ar}(n) = \mathrm{Conf}_n(\C^3)$ (unordered configurations of $n$ points in $\C^3$).
\item Composition by nested inclusion of polydisks $D_r(z) = \{w : |w_i - z_i| < r_i\}$.
\item $\mathbb{T}^3$-equivariance by the diagonal torus action on $\C^3$.
\end{itemize}
$\pi_1(\mathrm{Conf}_2(\C^3)) = \pi_1(S^5) = 0$ (Hurewicz), so $E_3^{\hol}$ is simply connected: no braid monodromy, but the Koszul $(-2)$-shifted bracket is nonzero. Over $\Q$, $H_*(E_3^{\hol}; \Q) \simeq H_*(E_3^{\top}; \Q) \simeq \Pois_3$, the shifted Poisson operad; over chain complexes this is not equivalent to commutative algebras.
\end{definition}

\begin{proposition}[Koszul dual of $E_3^{\hol}$]
\label{prop:e3-hol-koszul}
The Koszul dual operad of $E_3^{\hol}$ in chain complexes is $(E_3^{\hol})^! \simeq \Lie[2]$, the $2$-shifted Lie operad. In particular, an $E_3^{\hol}$-algebra in chain complexes is modelled by a chain complex $\obsc$ with a graded-commutative product and a Poisson bracket $\{-,-\}$ of degree $-2$ satisfying Leibniz. (Fresse 2017 Vol.~I Thm.~14.1.A.)
\end{proposition}

\begin{remark}[CFG vs avatar operad]
\label{rem:cfg-vs-avatar-operad}
CFG 2026 works with $E_3^{\top} = $ topological little $3$-disks on $\R^3$: same $\pi_1 = 0$, same rational Poisson-$3$ homotopy type, but its factorisation algebras on $\R^3$ recover all topological $E_3$-algebras. The avatar operad $E_3^{\hol}$ recovers only the Dolbeault-holomorphic sector on $\C^3$. Every theorem on the CFG side that depends only on operadic $E_3$-homotopy type transports across; theorems that depend on full topological locality (e.g.\ smooth-manifold invariance of factorisation homology) do not transport without additional holomorphic-locality hypotheses.
\end{remark}

\subsubsection{Classical BV datum from first principles}
\label{sssec:classical-BV-datum}

\begin{definition}[hCS BV complex on $\C^3$]
\label{def:hCS-BV-avatar}
Fix a finite-dimensional Lie algebra $\fg$ with invariant non-degenerate pairing $\langle -, - \rangle$ and holomorphic volume $\Omega = dz_1 \wedge dz_2 \wedge dz_3$. The classical BV field complex is
\[
\cE_{\hCS} := \Omega^{0,\bullet}(\C^3, \fg)[1],
\]
graded by BV ghost number: $c \in \Omega^{0,0}(\fg)[1]$ (ghost, gh $= -1$), $A \in \Omega^{0,1}(\fg)$ (gauge field, gh $= 0$), $A^* \in \Omega^{0,2}(\fg)[-1]$ (antifield, gh $= +1$), $c^* \in \Omega^{0,3}(\fg)[-2]$ (antighost, gh $= +2$). Pair via
\[
\omega(\alpha, \beta) := \int_{\C^3} \Omega \wedge \langle \alpha, \beta \rangle,
\]
a $(-1)$-shifted symplectic pairing by Serre duality $H^{0,3}(\C^3, \omega_c) \simeq \C$.
\end{definition}

\begin{proposition}[Classical master equation]
\label{prop:classical-ME-avatar}
The classical action
\[
S_{\mathrm{cl}}[\cA] \;=\; \int_{\C^3} \Omega \wedge \left( \tfrac{1}{2} \langle \cA, \bar\partial \cA \rangle + \tfrac{1}{6} \langle \cA, [\cA, \cA] \rangle \right), \qquad \cA \in \cE_{\hCS}
\]
satisfies $\{S_{\mathrm{cl}}, S_{\mathrm{cl}}\} = 0$ by the Jacobi identity on $\fg$ and integration by parts against $\bar\partial\Omega = 0$. Expanding $\cA = c + A + A^* + c^*$ and reading the $(3,3)$-form component produces the usual BRST structure: $Q_0 c = \tfrac{1}{2}[c,c]$, $Q_0 A = \bar\partial c + [c, A]$, $Q_0 A^* = [c, A^*] + \bar\partial A + \tfrac{1}{2}[A, A]$, $Q_0 c^* = [c, c^*] + \bar\partial A^* + [A, A^*]$.
\end{proposition}

\subsubsection{Maurer--Cartan, twisting, $L_\infty$-obstruction}
\label{sssec:MC-twisting-Linf}

\begin{proposition}[Maurer--Cartan = flat hCS]
\label{prop:MC-flat-hCS-avatar}
The classical BV field complex $\cE_{\hCS}$ is a dg $L_\infty$-algebra with:
\begin{itemize}
\item $\ell_1(\cA) = \bar\partial \cA$;
\item $\ell_2(\cA_1, \cA_2) = [\cA_1, \cA_2]$;
\item $\ell_k = 0$ for $k \geq 3$ (on $\C^3$; compact CY$_3$ acquires $\ell_3^{\min} = Y_3$ via Kapranov).
\end{itemize}
The Maurer--Cartan set
$\mathrm{MC}(\cE_{\hCS}) = \{\cA : \bar\partial \cA + \tfrac{1}{2}[\cA, \cA] = 0\}$
is the locus of solutions to the classical hCS equation of motion, and the gauge quotient
$\mathrm{MC}(\cE_{\hCS}) / \mathcal{G}$ is the moduli of holomorphic $\fg$-bundles on $\C^3$ with BV gauge-fixing.
\end{proposition}

\begin{proposition}[Twisting and formal deformation theory]
\label{prop:twisting-hCS-avatar}
For $\cA_0 \in \mathrm{MC}(\cE_{\hCS})$, the twisted $L_\infty$-algebra $(\cE_{\hCS}^{\cA_0}, \ell_1^{\cA_0}, \ell_2^{\cA_0}, \ldots)$ has
\[
\ell_1^{\cA_0}(\alpha) = \bar\partial \alpha + [\cA_0, \alpha], \qquad \ell_2^{\cA_0} = [\cdot, \cdot],
\]
and formal deformations of $\cA_0$ are classified by $H^1(\cE_{\hCS}^{\cA_0})$; obstructions live in $H^2(\cE_{\hCS}^{\cA_0})$. For $\cA_0 = 0$ on $\C^3$: $H^1 = H^{0,1}(\C^3, \fg) = 0$ (rigid), $H^2 = H^{0,2}(\C^3, \fg) = 0$ (unobstructed).
\end{proposition}

\subsubsection{The BV--BRST complex and factorisation algebra}
\label{sssec:BV-BRST-factorisation}

\begin{definition}[Classical observables as CE cochains]
\label{def:classical-obs-CE-avatar}
\[
\mathrm{Obs}^{\mathrm{cl}}_{\hCS}(U) := \mathrm{CE}^\bullet(\cE_{\hCS}(U)) = \widehat{\mathrm{Sym}}^\bullet\!\left(\cE_{\hCS}(U)^\vee[-1]\right), \quad U \subset \C^3 \text{ open},
\]
with Chevalley--Eilenberg differential $Q_0$ dual to the $L_\infty$-structure. This is the algebra of formal functions on the $L_\infty$-field complex, with $(-1)$-shifted symplectic structure induced from $\omega$.
\end{definition}

\begin{theorem}[Classical prefactorisation algebra]
\label{thm:classical-prefact-avatar}
$U \mapsto \mathrm{Obs}^{\mathrm{cl}}_{\hCS}(U)$ forms a holomorphic prefactorisation algebra on $\C^3$: for disjoint $U_1, \ldots, U_n \subset V$, the structure maps
\[
\mathrm{Obs}^{\mathrm{cl}}_{\hCS}(U_1) \otimes \cdots \otimes \mathrm{Obs}^{\mathrm{cl}}_{\hCS}(U_n) \to \mathrm{Obs}^{\mathrm{cl}}_{\hCS}(V)
\]
are the pullback along pullback of the $L_\infty$-field complex, compatible with $Q_0$ and holomorphic with respect to the Dolbeault complex structure. (Costello--Gwilliam 2017 Vol.~II §5.1.)
\end{theorem}

\subsubsection{Costello renormalisation in the chiral avatar}
\label{sssec:costello-renormalisation-avatar}

\begin{definition}[Heat-kernel regulator]
\label{def:heat-kernel-avatar}
For $\varepsilon > 0$ let $K_\varepsilon : \C^3 \times \C^3 \to \wedge^3 T^*_{\C^3} \otimes \wedge^3 T^*_{\C^3}$ be the heat kernel for the Laplacian $\Delta_{\bar\partial}$ at time $\varepsilon$. The regularised propagator is
\[
P_{\varepsilon, L}(z, w) = \int_\varepsilon^L (\bar\partial^* \otimes 1) K_t(z, w) \, dt,
\]
degenerating at $\varepsilon \to 0$ to the Bochner--Martinelli kernel
\[
P_{\mathrm{BM}}(z, w) = \frac{2}{(2\pi i)^3} \sum_{k=1}^{3} \frac{(-1)^{k-1} \overline{(z_k - w_k)}}{\|z - w\|^6} \, \widehat{d\bar z_k} \wedge dw_1 \wedge dw_2 \wedge dw_3.
\]
\end{definition}

\begin{theorem}[Renormalisation group flow and counterterms]
\label{thm:RG-flow-avatar}
Costello's RG flow $W(P_{\varepsilon, L}, S) = \hbar \log \int e^{(S + I)/\hbar}$ extends the classical action $S_{\mathrm{cl}}$ to a family of effective interactions $\{I[L]\}_{L>0}$ satisfying:
\begin{enumerate}
\item \emph{RG equation:} $I[L'] = W(P_{L, L'}, I[L])$.
\item \emph{Scale-zero locality:} $\lim_{L \to 0} I[L]$ exists modulo local counterterms $\hbar \sum_n \hbar^n I^{\mathrm{c.t.}}_n[\varepsilon]$.
\item \emph{Quantum master equation:} $(Q_0 + \hbar \Delta) e^{(S + I[L])/\hbar} = 0$, equivalently $Q_0 I + \tfrac{1}{2}\{I, I\} + \hbar \Delta I = 0$.
\end{enumerate}
The space of counterterms is parametrised by a finite-dimensional renormalisation group; on $\C^3$ with gauge $\fgl_r$ the one-loop counterterm is the wave-function $- \hbar C_2(\fg)(4\pi)^{-3} \log(L/\varepsilon) \int \Omega \wedge \mathrm{Tr}(\cA \bar\partial \cA)$ plus the BCOV one-loop curving $\chi(X)/24 \cdot [\Omega]^{0,1}$ in the compact case.
\end{theorem}

\subsubsection{The quantum $E_3$ holomorphic factorisation algebra}
\label{sssec:quantum-E3-factorisation}

\begin{theorem}[Quantum observables as $E_3$-holomorphic factorisation algebra]
\label{thm:quantum-obs-E3-avatar}
At each renormalisation scale $L$, the quantum observables
\[
\mathrm{Obs}^{\mathrm{q}}_{\hCS}(U)[L] := \left( \mathrm{CE}^\bullet(\cE_{\hCS}(U))[[\hbar]], \; Q_0 + \hbar \{I[L], -\} + \hbar \Delta_L \right)
\]
form an $E_3$-holomorphic factorisation algebra on $\C^3$ in the sense of Costello--Gwilliam Vol.~II §5.3: disjoint-union structure maps are chain maps, Dolbeault-holomorphic in $z$, and the operad action by $\mathrm{Conf}_n(\C^3) / \mathbb{T}^3$ is coherent through the Bochner--Martinelli Feynman weights.

\noindent
\emph{Proof sketch.} The prefactorisation structure is Theorem~\ref{thm:classical-prefact-avatar} transported to the quantum level via the Costello RG (Theorem~\ref{thm:RG-flow-avatar}); holomorphic locality is the Dolbeault compatibility of $P_{\mathrm{BM}}$ (Definition~\ref{def:heat-kernel-avatar}); the $E_3$-operad coherence is the integration-over-$\overline{\mathrm{Conf}}_n(\C^3)$ formula for the Feynman amplitudes, which agree with the Kontsevich graph-weight formula on $\R^3$ upon restriction to real configurations, and hence pick the Kontsevich associator in the $\mathrm{GRT}_1(\Q)$-torsor (Tamarkin 1998, Willwacher 2014).
\end{theorem}

\begin{corollary}[Generators and relations]
\label{cor:E3-gens-rels-avatar}
As an $E_3^{\hol}$-algebra in chain complexes, $\mathrm{Obs}^{\mathrm{q}}_{\hCS}(\C^3)$ is generated by the linear observables
\[
\cO_\phi[\cA] := \int_{\C^3} \Omega \wedge \langle \phi, \cA \rangle, \qquad \phi \in \Omega^{0,\bullet}_c(\C^3, \fg^\vee)
\]
subject to the relations:
\begin{enumerate}
\item $\cO_{\bar\partial \phi} = Q_0 \cO_\phi$ (compatibility with $\ell_1$);
\item $\cO_\phi \star \cO_\psi - (-1)^{|\phi||\psi|} \cO_\psi \star \cO_\phi = \hbar \cO_{[\phi, \psi]_{\mathrm{BM}}}$ (Poisson bracket via Bochner--Martinelli propagator) up to quantum corrections;
\item $\cO_\phi \star \cO_\psi \star \cO_\chi = \sum_\Gamma w_\Gamma \cdot B_\Gamma(\phi, \psi, \chi)$ with $w_\Gamma = \int_{\overline{\mathrm{Conf}}_n(\C^3)} \omega_\Gamma$ (the Kontsevich cubic).
\end{enumerate}
These exhibit $\mathrm{Obs}^{\mathrm{q}}_{\hCS}(\C^3)$ as the $E_3^{\hol}$-enveloping algebra $U^{E_3^{\hol}}(\cE_{\hCS}[1])$ of the $L_\infty$-algebra $\cE_{\hCS}$ with its shifted pairing.
\end{corollary}

\begin{theorem}[BV--BRST trace]
\label{thm:BV-BRST-trace-avatar}
For compact $X$ a Calabi--Yau threefold, the BV--BRST trace on $\mathrm{Obs}^{\mathrm{q}}_{\hCS}(X)$ is
\[
\mathrm{Tr}_{\mathrm{BV}}(\cO) = \int_{\mathcal{L}} e^{(S + I[L])/\hbar} \cdot \cO
\]
where $\mathcal{L} \subset \cE_{\hCS}(X)$ is a $(-1)$-shifted Lagrangian (gauge-fixing). On the Stage-$1$ factorisation algebra this agrees with the $E_3$-trace determined by the nondegenerate pairing $S^5 \subset \partial \overline{\mathrm{Conf}}_2(\C^3)$, giving a trace class $\mathrm{Tr}_{E_3} \in \mathrm{HH}_{-3}^{E_3}(\mathrm{Obs}^{\mathrm{q}}_{\hCS})$.
\end{theorem}

\subsubsection{Cross-consistency: CoHA, $\mathcal{W}_{1+\infty}$, and the $E_3$ avatar}
\label{sssec:cross-consistency-CoHA-W-E3}

\begin{theorem}[Avatar cross-consistency with CoHA chain]
\label{thm:cross-consistency-avatar}
The three structures
\[
\begin{array}{c}
\mathrm{CoHA}(\C^3) \;\simeq\; Y^+(\widehat{\fgl}_1) \\
\downarrow \text{\small(Drinfeld double)} \\
Y(\widehat{\fgl}_1) \;\simeq\; \mathcal{W}_{1+\infty} \text{ affine Yangian} \\
\updownarrow \text{\small(hCS-to-Hall comparison, conjectural)} \\
\mathrm{Obs}^{\mathrm{q}}_{\hCS}(\C^3) \text{ in } E_3^{\hol}\text{-}\HolFA(\C^3)
\end{array}
\]
are compatible in the following precise sense:
\begin{enumerate}
\item $\mathrm{CoHA}(\C^3) \simeq Y^+$ only; the full $\mathcal{W}_{1+\infty}$ requires the Drinfeld double, NOT $\mathrm{CoHA}$ directly (cache discipline).
\item $\mathrm{Obs}^{\mathrm{q}}_{\hCS}(\C^3)$ is an $E_3^{\hol}$-algebra on $\C^3$; $\mathcal{W}_{1+\infty}$ is an $E_2$-(vertex) algebra on $\C$. The passage $E_3^{\hol}(\C^3) \to E_2(\C)$ is Stage-$2$ factorisation homology along the transverse $\C^2$ fibre:
\[
\int_{\C^2} \mathrm{Obs}^{\mathrm{q}}_{\hCS}(\C^3) \;\simeq\; \mathrm{Obs}^{\mathrm{q}}_{\mathrm{defect}}(\C), \quad \text{a vertex algebra.}
\]
\item The identification $\mathrm{Obs}^{\mathrm{q}}_{\mathrm{defect}}(\C) \simeq \mathcal{W}_{1+\infty}$ is the Costello 2013 + Gaiotto--Rap\v{c}\'ak 2018 conjectural defect identification, not a direct theorem of CFG 2026.
\end{enumerate}
\end{theorem}

\begin{remark}[AP-CY discipline at $d = 3$]
\label{rem:AP-CY-d3-avatar}
At $d = 3$, the CY chiral algebra $A_{\C^3} = \Phi_3(\Perf(\C^3))$ is $E_1$. The $E_3^{\hol}$-structure on $\mathrm{Obs}^{\mathrm{q}}_{\hCS}(\C^3)$ is a statement about the \emph{factorisation algebra of the 6d hCS field theory on $\C^3$}, not about $A_{\C^3}$ itself. The $E_2$-braided enhancement, when it exists, lives on the Drinfeld centre $\cZ(\mathrm{Rep}^{E_1}(A_{\C^3}))$; at the field-theoretic level this is the statement that Stage-$2$ specialisation $\int_{\C^2}$ produces an $E_1 = E_2 \otimes E_{-1}$-structure, with the $E_2$-shadow recovered on $\cZ$. Conflating the ambient $E_3^{\hol}$-structure on the factorisation algebra with an $E_2$-structure on $A$ is AP-CY: the ambient field-theory operad is not the CY-category-output operad.
\end{remark}

\emph{End of treatise.}
